% =========================================================
% Proof: Nested SB Intervals Contain a Unique Limit Point
% =========================================================

\newpage
\phantomsection
\label{prf:sb-unique-point}

\begin{remark*}[Return]
\hyperref[prop:sb-unique-point]{Return to Proposition~\ref*{prop:sb-unique-point}.}
\end{remark*}

\begin{theorem*}[Nested Stern--Brocot intervals contain a unique limit point]
Let $(I_n)$ with $I_n = [a_n/b_n, c_n/d_n]$ be the nested closed rational
intervals from the SB refinement process.  Then $\bigcap_{n=0}^{\infty} I_n$
contains exactly one real number.
\end{theorem*}

\begin{proof}[Proof (Professional Standard)]
\textbf{Existence.} The sequence $(a_n/b_n)$ is monotone increasing and
bounded above by $c_0/d_0$; the sequence $(c_n/d_n)$ is monotone decreasing
and bounded below by $a_0/b_0$.  By the Monotone Convergence Theorem in
$\mathbb{R}$, both sequences converge:
\[
\alpha = \lim_{n\to\infty} \frac{a_n}{b_n} \in \mathbb{R},
\qquad
\beta = \lim_{n\to\infty} \frac{c_n}{d_n} \in \mathbb{R}.
\]
Since $a_n/b_n \leq c_n/d_n$ for all $n$, we have $\alpha \leq \beta$.
But
\[
0 \leq \beta - \alpha \leq \ell_n = \frac{1}{b_n d_n}
\]
for all $n$, and $\ell_n \to 0$
by \hyperref[lem:sb-length-zero]{Lemma~\ref*{lem:sb-length-zero}}.
Therefore $\alpha = \beta$.  Set $\xi = \alpha = \beta$; then
$a_n/b_n \leq \xi \leq c_n/d_n$ for all $n$, so $\xi \in \bigcap I_n$.

\textbf{Uniqueness.} If $\xi, \xi' \in \bigcap I_n$, then for every $n$:
\[
|\xi - \xi'| \leq \ell_n \to 0.
\]
Therefore $\xi = \xi'$.
\end{proof}

\begin{proof}[Proof (Detailed Learning)]
\textbf{Step 1: Show the left endpoints form a bounded increasing sequence.}
By the Stern--Brocot construction rule, the left endpoint is either
unchanged or replaced by the mediant, which lies strictly to the right of
the previous left endpoint (by the Mediant Inequality).  So $(a_n/b_n)$
is non-decreasing.  It is bounded above by $c_0/d_0$ since $a_n/b_n <
c_n/d_n \leq c_0/d_0$ at every step.

\textbf{Step 2: Show the right endpoints form a bounded decreasing sequence.}
Symmetrically, $(c_n/d_n)$ is non-increasing and bounded below by $a_0/b_0$.

\textbf{Step 3: Both sequences converge in $\mathbb{R}$.}
By the Monotone Convergence Theorem in $\mathbb{R}$ (which requires the
completeness of $\mathbb{R}$):
\[
\alpha = \sup_n \frac{a_n}{b_n} \in \mathbb{R}, \qquad
\beta = \inf_n \frac{c_n}{d_n} \in \mathbb{R}.
\]

\textbf{Step 4: The two limits coincide.}
For every $n$: $\alpha \leq \beta$ (since $a_n/b_n \leq c_n/d_n$ for
all $n$, and $\alpha, \beta$ are the respective limiting extremes).
Also:
\[
0 \leq \beta - \alpha \leq \frac{c_n}{d_n} - \frac{a_n}{b_n} = \ell_n.
\]
Since $\ell_n \to 0$ by \hyperref[lem:sb-length-zero]{Lemma~\ref*{lem:sb-length-zero}},
the squeeze theorem gives $\beta - \alpha = 0$, so $\alpha = \beta$.
Set $\xi = \alpha = \beta$.

\textbf{Step 5: Establish $\xi \in \bigcap I_n$.}
For every $n$: $a_n/b_n \leq \xi \leq c_n/d_n$, so $\xi \in I_n$.
Therefore $\xi \in \bigcap_{n=0}^{\infty} I_n$.

\textbf{Step 6: Uniqueness.}
Suppose $\xi, \xi' \in \bigcap I_n$.  For every $n$, both lie in $I_n$,
so $|\xi - \xi'| \leq \ell_n$.  Since $\ell_n \to 0$, we have
$|\xi - \xi'| = 0$, hence $\xi = \xi'$.
\end{proof}

\begin{remark*}[Proof structure]
A nested intervals argument in $\mathbb{R}$: monotone bounded sequences
converge, the two limits must coincide because the gaps shrink to zero,
and uniqueness is a direct squeeze.  The completeness of $\mathbb{R}$
(via the Monotone Convergence Theorem) is essential — this argument
fails in $\mathbb{Q}$ because $\xi$ need not be rational.
\end{remark*}

\begin{remark*}[Dependencies]
Monotone Convergence Theorem in $\mathbb{R}$ (requires completeness of $\mathbb{R}$),
\hyperref[lem:sb-length-zero]{Stern--Brocot interval lengths shrink to zero},
\hyperref[thm:mediant-inequality]{Mediant inequality}.
\end{remark*}
